% TOP_PROBLEM: {"id": 3028, "title": "Finite Lattice Representation Problem", "category_id": 4, "category": {"id": 4, "name": "algebra", "display_name": "Algebra", "description": "Group theory, ring theory, field theory, and algebraic structures.", "slug": "algebra", "order_index": 4, "created_at": "2026-07-31T15:26:25.671Z"}, "status": "open", "problem_number": "OPG-37283", "set_id": 12}
% FIRST_POSTED: 2026-10-01
% ATTEMPT_STATUS: unresolved
% UnsolvedMath ID 3028; source catalogue code OPG-37283
% !TeX program = lualatex
% Research attempt by GPT-6 Astra Ultra; no claim of a new complete solution.
\documentclass[11pt,a4paper]{article}
\usepackage[margin=25mm]{geometry}
\usepackage{fontspec}
\setmainfont{lmroman10-regular.otf}[BoldFont=lmroman10-bold.otf,
 ItalicFont=lmroman10-italic.otf,BoldItalicFont=lmroman10-bolditalic.otf]
\usepackage{amsmath,amssymb,mathtools}
\usepackage{unicode-math}
\setmathfont{latinmodern-math.otf}
\AtBeginDocument{\let\setminus\smallsetminus}
\usepackage{xurl,hyperref}
\hypersetup{colorlinks=true,urlcolor=blue}
\setcounter{secnumdepth}{0}
\setlength{\parindent}{0pt}
\setlength{\parskip}{5pt}
\setlength{\emergencystretch}{3em}
\begin{document}
\section*{Finite Lattice Representation Problem}\label{3028}
{\small UnsolvedMath ID 3028 · OPG-37283 · Algebra · OpenGarden}
\subsection{Definitions and mathematical statement}
A finite lattice is a finite nonempty partially ordered set $L$ in which every
pair has a greatest lower bound $x\wedge y$ and a least upper bound $x\vee y$.
An algebra $\mathbf A=(A;(f_i)_{i\in I})$ consists of a nonempty set $A$ together
with specified operations $f_i:A^{n_i}\to A$ of finite arities $n_i$.
Here ``finite algebra'' means that $A$ is finite; its operations may be chosen
as part of a proposed representation.

A congruence of $\mathbf A$ is an equivalence relation $\theta$ on $A$ compatible
with every operation: whenever $a_j\mathrel\theta b_j$ for all inputs, one has
$f_i(a_1,\ldots,a_{n_i})\mathrel\theta f_i(b_1,\ldots,b_{n_i})$.
The set $\operatorname{Con}(\mathbf A)$, ordered by inclusion, is a lattice.
Its meet is intersection, and its join is the smallest congruence containing
both relations. A representation requires an isomorphism of lattices, hence
preservation of both meet and join, and not merely an embedding.

The catalogue conjecture is the negative form of the finite lattice
representation problem:
\[
 \exists\text{ finite lattice }L\quad
 \forall\text{ finite algebras }\mathbf A,\qquad
 L\not\cong\operatorname{Con}(\mathbf A).
\]
Thus a positive resolution of this conjecture must exhibit a finite lattice
and exclude every possible finite underlying set and collection of operations.
Its negation says that every finite lattice admits such a representation;
the size of the representing algebra is not bounded in advance. No fixed
algebraic signature, such as groups or rings alone, is imposed.
\subsection{Short English statement}
Is there a finite lattice that cannot be the lattice of compatible equivalence relations of any finite algebra?
\subsection{Research attempt}
\paragraph{Scope and process.}
This is a GPT-6 Astra Ultra research attempt dated 1 October 2026, using
primary-source browsing and elementary mathematical checks. The canonical
definition above is retained. Results below are restricted results or
reductions, not a claim of a new solution to the entire catalogue question.
No formal proof assistant or external mathematical referee checked this text.


\paragraph{Literature and target.}
The unrestricted finite-algebra question is the one in [S2]. DeMeo's
research monograph [R1] discusses its relationship with subgroup intervals;
Chan [R2] explicitly records representability of every finite distributive
lattice. We give a direct construction of that known restricted result,
with a bound on the size of the underlying set. No assertion about a
particular smallest unknown lattice is needed here.

\paragraph{Plan.}
A failed search within groups alone cannot exclude arbitrary algebras.
Instead, choose operations to force the congruences to be coordinate
subspaces, then encode an arbitrary finite partial order by additional
operations. This tests which obstruction a proposed negative example
would have to escape.

\paragraph{A finite algebra for any lattice of order ideals.}
Let $P$ be a finite partially ordered set and let $V=\mathbb F_2^P$.
Equip $V$ with addition, the constant zero, every coordinate projection
$\pi_p(x)=x_pe_p$, and, for every $q\le p$, the unary map
$T_{qp}(x)=x_pe_q$. These are finitely many operations on $2^{|P|}$
elements. For a congruence $\theta$, put $W=\{x:x\mathrel\theta0\}$.
Compatibility with addition shows that $W$ is a subspace and
\[
 x\mathrel\theta y\quad\Longleftrightarrow\quad x+y\in W.
\]
Indeed, add $y$ to both sides for one implication and add $y$ back for
the converse. Conversely, a subspace stable under all the stated linear
maps defines a compatible equivalence relation by this formula.

Projection stability implies
$W=\operatorname{span}\{e_p:p\in I\}$ for some $I\subseteq P$.
To see this, every coordinate present in a vector of $W$ can be isolated
by $\pi_p$, and every vector is the sum of its coordinate projections.
The maps $T_{qp}$ then imply $p\in I,\ q\le p\Rightarrow q\in I$.
Conversely any such down-set is stable under all operations. Thus
\[
 \operatorname{Con}(V;+,0,\pi_p,T_{qp})\cong J(P),
\]
where $J(P)$ is the lattice of down-sets. Inclusion, intersection and
union correspond exactly; this is an isomorphism, not just an embedding.

\paragraph{Recovering all finite distributive lattices.}
For completeness, let $L$ be finite distributive and let $P$ be its
nonzero join-irreducible elements. Every $x\in L$ is the join of the
members of $P$ below it, by induction on the finite order. Moreover a
join-irreducible $p$ is join-prime: if $p\le a\vee b$, then
$p=(p\wedge a)\vee(p\wedge b)$, so irreducibility gives $p\le a$ or
$p\le b$. Consequently
\[
 x\longmapsto\{p\in P:p\le x\}
\]
preserves meets and joins and is injective. A down-set $I$ is its image
at $x=\bigvee I$: join-primality ensures that no extra join-irreducible
lies below this join. This proves surjectivity, including the empty
down-set. Composing with the construction proves that $L$ has a finite
representation on at most $2^{|P|}$ elements.

\paragraph{Verification and obstruction to extending the argument.}
If $P$ is an antichain, only coordinate restrictions remain, yielding a
Boolean lattice. If $P$ is a chain, the possible subspaces are its initial
coordinate spans, yielding a chain of congruences. For $P$ empty the
one-element algebra gives the one-element lattice. These checks include
the extreme cases in the statement. The crucial join-primality step
uses distributivity: in the five-element diamond a join-irreducible
atom lies below the join of the other two atoms without lying below
either one. Our operations always give a distributive ideal lattice,
so this particular construction cannot represent that diamond. This is
an obstruction to the construction, not to all finite algebras.

\paragraph{Final status and first remaining gap.}
\textbf{UNRESOLVED.} Every finite distributive candidate is excluded
from being a counterexample by an explicit finite representation.
There is no argument here excluding a nondistributive lattice from all
finite algebras, nor a representation theorem for every nondistributive
lattice. Those are the missing alternatives required by the target.

\subsection{Sources}
\begin{itemize}
\item[S1] ULAM.AI and the UnsolvedMath Contributors, supplied problems.json, record 3028 (OPG-37283), OpenGarden collection. Catalogue identity and source transcription; CC BY 4.0.
\url{https://huggingface.co/datasets/ulamai/UnsolvedMath}.
\item[S2] Open Problem Garden, Finite Lattice Representation Problem. Original problem and discussion; the mathematical formulation below is expanded from the supplied source record.
\url{http://www.openproblemgarden.org/op/finite_congruence_lattice_problem}.

\item[R1] William J. DeMeo, \emph{Congruence lattices of finite algebras},
2012. Primary research monograph and context; abstract checked 1 October 2026.
\url{https://arxiv.org/abs/1204.4305}.
\item[R2] Brian T. Chan, \emph{Congruence Lattices of Certain Finite
Algebras with Three Commutative Binary Operations}, 2014. Primary
source confirming the known distributive subcase.
\url{https://arxiv.org/abs/1409.5920}.

\end{itemize}
\end{document}
